% Part: first-order-logic
% Chapter: model-theory
% Section: partial-iso

\documentclass[../../../include/open-logic-section]{subfiles}

\begin{document}

\olfileid{mod}{bas}{pis}
\section{Isomorfisme Parsial}

\begin{defn}
  Yen diwenehi rong !!{structure}s $\Struct{M}$
  lan $\Struct{N}$ kanggo basa kang padha, \emph{isomorfisme parsial}
  saka $\Struct{M}$ menyang $\Struct{N}$
  iku fungsi parsial winates $p$
  kang nampa argumen saka $\Domain M$
  lan ngasilake nilai ing $\Domain N$,
  lan nyukupi syarat-syarat isomorfisme
  saka \olref[iso]{defn:isomorphism}
  ing ranahé:
  \begin{enumerate}
  \item $p$ iku !!{injective};
  \item kanggo saben !!{constant}~$c$:
    yen $p(\Assign{c}{M})$ katetepake,
    mula $p(\Assign{c}{M}) = \Assign{c}{N}$;
  \item kanggo saben $n$-panggonan
    !!{predicate} $P$: yen $a_1$, \dots, $a_n$
    ana ing ranah $p$, mula
    $\langle a_1, \dots, a_n\rangle \in
    \Assign P M$ yen lan mung yen
    $\langle p(a_1), \dots, p(a_n) \rangle
    \in \Assign P N$;
  \item kanggo saben $n$-panggonan
    !!{function} $f$: yen $a_1$, \dots, $a_n$
    ana ing ranah $p$, mula
    $p(\Assign f M (a_1, \dots,a_n))
    = \Assign f N (p(a_1), \dots, p(a_n))$.
  \end{enumerate}
  $p$ winates tegesé $\dom{p}$ winates. Cathetan cakupan: miturut klausa fungsi ing definisi iki, ranah p kudu katutup tumrap fungsi kang diterapake marang tuple saka ranahe. Iki syarat kang luwih mbatesi kanggo fungsi kang duwe orbit tanpa wates; definisi alternatif kang nggunakake isomorfisme antarane substruktur kang dibangkitake ora diganti kanthi meneng ing edhisi iki.
\end{defn}

Gatekna yen fungsi kosong~$\emptyset$
mesthi dadi isomorfisme parsial
antarane sembarang rong !!{structure}s.

\begin{defn}\ollabel{defn:partialisom}
  Rong !!{structure}s $\Struct{M}$
  lan $\Struct{N}$ diarani
  \emph{isomorfis parsial}, ditulis
  $\Struct{M} \iso[p]
  \Struct{N}$, yen lan mung yen
  ana himpunan ora kosong $I$
  kang isiné isomorfisme parsial
  antarane $\Struct{M}$ lan $\Struct{N}$
  lan nyukupi sipat \emph{maju-mundur}:
  \begin{enumerate}
  \item (\emph{Maju}) Kanggo saben $p \in I$
    lan $a \in \Domain M$,
    ana $q \in I$ saéngga
    $p \subseteq q$ lan $a$
    ana ing ranah $q$;
  \item (\emph{Mundur}) Kanggo saben $p \in I$
    lan $b \in \Domain N$,
    ana $q \in I$ saéngga
    $p \subseteq q$ lan $b$
    ana ing jajaran asil $q$.
  \end{enumerate}
\end{defn}

\begin{thm}\ollabel{thm:p-isom1}
  Yen $\Struct{M} \iso[p] \Struct{N}$
  lan $\Struct{M}$ sarta
  $\Struct{N}$ padha !!{enumerable},
  mula $\Struct{M} \iso
  \Struct{N}$.
\end{thm}

\begin{proof}
  Amarga $\Struct{M}$ lan
  $\Struct{N}$ padha !!{enumerable},
  jupuken $\Domain{M} =
  \{a_0, a_1, \ldots \}$
  lan $\Domain{N} = \{b_0, b_1, \ldots \}$.
  Wiwit saka sembarang $p_0 \in I$,
  kita netepake urutan isomorfisme parsial
  kang saya amba $p_0 \subseteq p_1
  \subseteq p_2 \subseteq \cdots$
  kaya mangkene:
  \begin{enumerate}
  \item yen $n+1$ ganjil, upamane $n = 2r$,
    kanthi sipat Maju goleka
    $p_{n+1} \in I$ saéngga
    $p_n \subseteq p_{n+1}$
    lan $a_r$ ana ing ranah $p_{n+1}$;
  \item yen $n+1$ genep, upamane
    $n+1 =2r$, kanthi sipat Mundur
    goleka $p_{n+1} \in I$ saéngga
    $p_n \subseteq p_{n+1}$
    lan $b_{r-1}$ ana ing jajaran
    asil $p_{n+1}$.
  \end{enumerate}
  Cathetan koreksi sumber OLPL-089:
  nalika langkah genep kapisan,
  unsur kapisan saka jajaran asil
  kudu dilebokake; indeks ing
  sumber diwiwiti saka unsur kapindho.
Yen saiki ditetepake:
\[
p = \bigcup_{n\ge 0} p_n,
\]
mula $p$ iku isomorfisme
antarane $\Struct{M}$
lan $\Struct{N}$.
\end{proof}

\begin{prob}
  Tuduhna kanthi rinci yen $p$
  kang ditetepake ing
  \olref[mod][bas][pis]{thm:p-isom1}
  pancen isomorfisme.
\end{prob}

\begin{thm}\ollabel{thm:p-isom2}
  Upamakna $\Struct{M}$ lan
  $\Struct{N}$ iku !!{structure}s
  kanggo !!{language} relasional
  murni (!!a{language} kang mung
  ngemot !!{predicate}s,
  tanpa !!{function}s utawa konstanta).
  Yen $\Struct{M} \iso[p] \Struct{N}$,
  mula uga $\Struct{M} \elemequiv
  \Struct{N}$.
\end{thm}

\begin{proof}
  Kanthi induksi tumrap !!{formula}s,
  kita nuduhake yen $a_1$, \dots, $a_n$
  lan $b_1$, \dots, $b_n$
  nduweni isomorfisme parsial $p \in I$
  kang ngirim saben $a_i$ menyang $b_i$,
  lan $s_1(x_i) =a_i$ sarta
  $s_2(x_i) =b_i$ (kanggo $i =1$,
  \dots,~$n$), lan kabeh variabel bebas saka formula kang dibandhingake ana ing dhaptar kasebut, mula
  $\Sat{M}{!A}[s_1]$ yen lan
  mung yen $\Sat{N}{!A}[s_2]$.
  Cathetan cakupan SOL6-F018: p kudu anggota kulawarga maju-mundur kang ana ing hipotesis teorema. Isomorfisme parsial sembarang mung njaga data atomik ing ranahe; kanggo langkah kuantor, p kudu bisa dijembarake ing kulawarga kasebut. Upamane, rong unsur kang padha ora nyukupi predikat bisa dicocogake dening siji p, senajan mung salah siji struktur nduweni unsur liyane kang nyukupi predikat kasebut. Formula eksistensial mula ora mesthi padha kanggo p sembarang.
  Kasus $n=0$ menehi
  $\Struct{M} \elemequiv \Struct{N}$.
\end{proof}

\begin{rem}
Yen ana !!{function}s, asil sadurunge
isih bener, nanging kita kudu
nggatekake isomorfisme kang diasilake
dening ekstensi $p$ saka kulawarga maju-mundur antarane
sub!!{structure} saka $\Struct{M}$
kang dibangkitake dening $a_1$, \dots, $a_n$
lan sub!!{structure} saka
$\Struct{N}$ kang dibangkitake
dening $b_1$, \dots, $b_n$.
\end{rem}

Asil sadurunge bisa dipérang
dadi sawetara tataran kanthi
ngubungake cacah kuantor
kang nyarang ing !!a{formula}
karo cacah ekstensi
isomorfisme parsial kang dibutuhake.

\begin{defn}
  Kanggo sembarang !!{formula}~$!A$,
  \emph{tingkat kuantor} saka $!A$,
  ditulis $\QuantRank{!A} \in \Nat$,
  ditetepake sacara rekursif minangka
  cacah paling gedhé kuantor kang
  nyarang ing $!A$. Rong
  !!{structure}s $\Struct{M}$
  lan $\Struct{N}$ diarani
  \emph{ekuivalen-$n$}, ditulis
  $\Struct{M} \elemequiv[n] \Struct{N}$,
  yen padha kanggo kabeh ukara
  kanthi tingkat kuantor
  kurang saka utawa padha karo~$n$.
\end{defn}

\begin{prop}\ollabel{prop:qr-finite}
  Upamakna $\Lang{L}$ iku
  !!{language} winates tanpa
  simbol fungsi, yaiku
  !!{language} kang ngemot
  !!{predicate}s lan !!{constant}s
  winates cacahé, lan tanpa
  !!{function}s. Mula kanggo
  saben $n \in \Nat$ mung ana
  winates cacahé !!{sentence}s
  tataran kapisan ing !!{language}
  $\Lang{L}$ kang tingkat
  kuantoré ora luwih saka $n$,
  nganti ekuivalensi logis.
  Cathetan koreksi sumber OLPL-090:
  ing proposisi iki konstanta diidinake,
  dadi teges basane luwih jembar
  tinimbang teges relasional murni
  ing teorema sadurunge.
\end{prop}

\begin{proof}
  Kanthi induksi tumrap $n$.
\end{proof}

\begin{defn}
  Yen diwenehi !!a{structure}~$\Struct{M}$,
  upamakna $\Domain M^{<\omega}$
  iku himpunan kabeh urutan winates
  saka $\Domain{M}$. Kita nggunakake
  $\mathbf{a}, \mathbf{b},
  \mathbf{c}, \ldots$
  kanggo nglambangake urutan winates
  saka unsur-unsur. Yen
  $\mathbf{a} \in \Domain{M}^{<\omega}$
  lan $a \in \Domain{M}$,
  mula $\mathbf{a}a$ nglambangake
  \emph{gabungan urutan}
  $\mathbf{a}$ karo $a$.
\end{defn}

\begin{defn}
  Yen diwenehi !!{structure}s
  $\Struct{M}$ lan $\Struct{N}$,
  kita netepake relasi
  $I_n \subseteq \Domain M^{<\omega}
  \times \Domain N^{<\omega}$
  antarane urutan-urutan kang padha
  dawane, kanthi rekursi tumrap $n$
  kaya mangkene:
   \begin{enumerate}
   \item $I_0(\mathbf{a},\mathbf{b})$
     yen lan mung yen $\mathbf{a}$
     lan $\mathbf{b}$ nyukupi
     !!{formula}s atomik kang padha
     ing $\Struct{M}$ lan
     $\Struct{N}$; yaiku, yen
     $s_1(x_i) = a_i$ lan
     $s_2(x_i) = b_i$
     sarta $!A$ iku atomik
     kanthi kabeh !!{variable}s
     ana ing antarane
     $x_1$, \dots,~$x_k$,
     mula $\Sat{M}{!A}[s_1]$
     yen lan mung yen~
     $\Sat{N}{!A}[s_2]$.
     Cathetan koreksi sumber OLPL-091:
     k ing indeks pungkasan iku
     dawa urutan, dudu tingkat rekursi.
   \item $I_{n+1} (\mathbf{a},\mathbf{b})$
     yen lan mung yen kanggo
     saben $a\in \Domain M$
     ana $b\in \Domain N$
     saéngga $I_n
     (\mathbf{a}a,\mathbf{b}b)$,
     lan uga kosok baline.
   \end{enumerate}
\end{defn}

\begin{defn}
  Tulis $\Struct{M} \approx_n
  \Struct{N}$ yen
  $I_n(\emptyseq,\emptyseq)$
  bener kanggo $\Struct{M}$
  lan $\Struct{N}$
  (ing kene $\emptyseq$
  iku urutan kosong).
\end{defn}

\begin{thm}\ollabel{thm:b-n-f}
  Upamakna $\Lang{L}$
  !!{language} relasional murni.
  Mula $I_n
  (\mathbf{a},\mathbf{b})$
  njalari yen kanggo saben $!A$
  kanthi $\QuantRank{!A} \le n$ lan saben variabel bebasé ana ing urutan kang dicocogake,
  $\Sat{M}{!A}[\mathbf{a}]$
  yen lan mung yen
  $\Sat{N}{!A}[\mathbf{b}]$
  (ing kene maneh $\mathbf{a}$
  nyukupi $!A$ yen sembarang
  $s$ kanthi
  $s(x_i) = a_i$
  kang nyukupi $!A$).
  Sabanjure, yen $\Lang{L}$
  winates, pratelan kosok
  baline uga bener. Cathetan koreksi sumber OLPL-754: pratelan iki mung kanggo formula kang variabel bebasé ana ing urutan kasebut; variabel bebas ing njabane ora dikontrol dening relasi kasebut.
\end{thm}

\begin{proof}
  Bukti yen $I_n(\mathbf{a},\mathbf{b})$
  njalari $\mathbf{a}$
  lan $\mathbf{b}$ nyukupi
  !!{formula}s kang padha
  kanthi tingkat kuantor
  ora luwih saka $n$
  lumaku kanthi induksi
  kang prasaja tumrap $!A$.
  Kanggo kosok baline,
  kita nggunakake induksi
  tumrap $n$ lan
  \olref{prop:qr-finite},
  kang njamin yen kanggo
  saben $n$ mung ana
  winates cacahé !!{formula}s
  kang ora ekuivalen
  kanthi tingkat kuantor kasebut, kanggo dhaptar variabel bebas kang tetep lan winates. Cathetan cakupan SOL6-F019: watesan iki perlu; kabeh formula kang variabel bebase ora diwatesi ora nduweni winates kelas ekuivalensi, malah ing tingkat nol. Proposisi ngenani ukara bisa ditrapake ing kene kanthi ngganti saben variabel bebas ing dhaptar tetep nganggo konstanta anyar; basane tetep winates lan tanpa fungsi.

  Kanggo $n=0$, hipotesis yen
  $\mathbf{a}$ lan $\mathbf{b}$
  nyukupi !!{formula}s
  tanpa kuantor kang padha
  njalari padha nyukupi
  formula atomik kang padha;
  mula $I_0(\mathbf{a},\mathbf{b})$.

  Kanggo kasus $n+1$,
  upamakna $\mathbf{a}$
  lan $\mathbf{b}$ nyukupi
  !!{formula}s kang padha
  kanthi tingkat kuantor
  ora luwih saka $n+1$.
  Kanggo nuduhake
  $I_{n+1}(\mathbf{a},\mathbf{b})$,
  cukup nuduhake yen kanggo
  saben $a \in \Domain M$
  ana $b \in \Domain N$
  saéngga $I_n(\mathbf{a}a,\mathbf{b}b)$.
  Miturut hipotesis induksi,
  iki padha karo nuduhake
  yen kanggo saben
  $a \in \Domain M$
  ana $b \in \Domain N$
  saéngga $\mathbf{a}a$
  lan $\mathbf{b}b$
  nyukupi !!{formula}s
  kang padha kanthi
  tingkat kuantor
  ora luwih saka $n$.

  Yen diwenehi $a \in \Domain M$,
  upamakna $!T^a_n$
  iku himpunan !!{formula}s
  $!B(x,\mathbf{y})$
  kanthi tingkat ora luwih saka $n$
  kang disukupi dening
  $\mathbf{a}a$ ing $\Struct{M}$.
  $!T^a_n$ mung nduweni
  winates kelas ekuivalensi logis;
  kita bisa milih wakil
  winates cacahé saka
  himpunan iki lan nganggep
  konjungsine minangka siji
  !!{formula} tataran kapisan.
  Cathetan koreksi sumber OLPL-092:
  himpunan formula sintaktis kasebut
  ora winates; kang winates
  yaiku kelas-kelas ekuivalensine.
  Mula $\mathbf{a}$ nyukupi
  $\lexists[x][!T^a_n(x,\mathbf{y})]$,
  kang tingkat kuantoré
  ora luwih saka $n+1$.
  Miturut hipotesis,
  $\mathbf{b}$ nyukupi
  !!{formula} kang padha
  ing $\Struct{N}$, mula
  ana $b \in \Domain N$
  saéngga $\mathbf{b}b$
  nyukupi $!T^a_n$;
  mliginé, $\mathbf{b}b$
  nyukupi !!{formula}s
  kang padha kanthi
  tingkat kuantor
  ora luwih saka $n$
  karo $\mathbf{a}a$.
  Kanthi cara padha,
  kanggo saben
  $b \in \Domain N$
  ana $a\in \Domain M$
  saéngga $\mathbf{a}a$
  lan $\mathbf{b}b$
  nyukupi !!{formula}s
  kang padha kanthi
  tingkat kuantor
  ora luwih saka $n$;
  mangkono bukti rampung.
\end{proof}

\begin{cor}\ollabel{cor:b-n-f}
  Yen $\Struct{M}$
  lan $\Struct{N}$
  iku !!{structure}s relasional
  murni ing !!{language} winates,
  mula $\Struct{M} \approx_n\Struct{N}$
  yen lan mung yen
  $\Struct{M} \elemequiv[n] \Struct{N}$.
  Mliginé $\Struct{M} \elemequiv
  \Struct{N}$ yen lan mung yen
  kanggo saben $n$,
  $\Struct{M} \approx_n \Struct{N}$.
\end{cor}

\end{document}
